\scanpage{51}
\mspage{\handwritten{- 13 -}}
% archive sheet number: 50
\[
\begin{aligned}
\handwritten{f}(u_1+w_1,\ldotp\ldotp,u_p+w_p)={}&\handwritten{f}(u_1,\ldotp\ldotp,u_p)+\handwritten{\sum}_{\handwritten{k=1}}^{p}D_k\handwritten{f}(0,\ldotp\ldotp,0)w_k+\handwritten{\sum}_{\handwritten{k=1}}^{p}\handwritten{f}_k(u_1,\ldotp\ldotp,u_p)w_k+\\
&+\handwritten{\sum}_{\handwritten{(n_j)(m_j)}}\handwritten{b}(n_1,\ldotp\ldotp,n_p,m_1,\ldotp\ldotp,m_p)u_1^{n_1}\ldotp\ldotp u_p^{n_p}w_1^{m_1}:\ldotp\ldotp w_p^{m_p}
\end{aligned}
\]
où les séries entières (absolument sommables) $\handwritten{f}_k$ n'ont pas de terme constant et où, dans la dernière série (absolument sommable) on a $\handwritten{\sum}_{\handwritten{k=1}}^{p}m_k\handwritten{\geq}2$ pour tous les termes. Supposons les séries absolument sommables pour $\handwritten{|}u_k\handwritten{|}=\handwritten{|}w_k\handwritten{|}=r$ $(1\handwritten{\leq}k\handwritten{\leq}p)$, et soit $M$ la borne supérieure des termes $\handwritten{\|}\handwritten{b}(n_1,\ldotp\ldotp,n_p,m_1,\ldotp\ldotp,m_p)\handwritten{\|}r^{\handwritten{\sum}_{k}(n_k+m_k)}$. Alors pour $\handwritten{|}u_k\handwritten{|}\handwritten{\leq}\handwritten{\frac{1}{2}}r$ et $\handwritten{|}w_k\handwritten{|}\handwritten{\leq}\handwritten{\frac{1}{2}}r$ $(1\handwritten{\leq}k\handwritten{\leq}p)$ on a, en posant $t=\sup_{\underset{\handwritten{\smile}}{k}}\handwritten{|}w_k\handwritten{|}$
\[
\begin{aligned}
&\handwritten{\sum}_{\handwritten{(n_j)(m_j)}}\handwritten{\|}\handwritten{b}(n_1,\ldotp\ldotp,n_p,m_1,\ldotp\ldotp,m_p)\handwritten{\|}\handwritten{\bigl|}u_1^{n_1}\ldots u_p^{n_p}w_1^{m_1}\ldotp\ldotp w_p^{m_p}\handwritten{\bigr|}\handwritten{\leq}\\
&\handwritten{\leq}Mt^2\handwritten{\sum}_{\handwritten{(n_j)(m_j)}}\left(\handwritten{\frac{1}{2}}\right)^{\handwritten{\sum}_{k}n_k}\left(\handwritten{\frac{t}{r}}\right)^{(\handwritten{\sum}_{k}m_k)-2}
\end{aligned}
\]
Comme la série qui figure au second membre de cette inégalité est convergente pour $\handwritten{|}t\handwritten{|}\handwritten{\leq}\handwritten{\frac{1}{2}}r$, on voit que la formule (6) résulte aussitôt de la continuité des séries entières.

COROLLAIRE.- \uline{Toute fonction holomorphe dans} $A$ \uline{est indéfiniment différentiable dans} $A$.

En effet, la prop.8 montre que toute fonction holomorphe $\handwritten{f}$ est différentiable en tout point, et que sa différentielle est l'application
\[
(u_1,\ldots,u_p)\handwritten{\longrightarrow}\handwritten{\sum}_{\handwritten{k=1}}^{p}D_k\handwritten{f}\,.u_k
\]
qu'on peut identifier à l'élément $(D_k\handwritten{f}(z_1,\ldotp\ldotp,z_p))_{1\handwritten{\leq}k\handwritten{\leq}p}$ de l'\uline{espace produit} ${E}^{p}$ en tout point. Or, il résulte de la prop.7 que cet élément est fonction holomorphe dans $A$, ce qui prouve que $\handwritten{f}$ est indéfiniment différentiable dans $A$.

Pour $p=1$, on dit qu'un point $a\handwritten{\in}K$ est un \uline{zéro d'ordre $h$} d'une fonction $\handwritten{f}$ holomorphe au voisinage de $a$ si, dans le développement de $\handwritten{f}$ en série entière suivant les puissances de $z-a$, le premier terme non nul a un exposant égal à $h$. Si $K$ est de caractéristique $0$, il est
